\documentclass[10pt,twocolumn]{article}
\usepackage[utf8]{inputenc}
\usepackage[margin=0.75in]{geometry}
\usepackage{amsmath,amssymb,amsfonts}
\usepackage{booktabs}
\usepackage{hyperref}
\usepackage{microtype}
\usepackage{cite}

\hypersetup{
    colorlinks=true,
    linkcolor=blue,
    citecolor=blue,
    urlcolor=blue
}

\title{\textbf{Combinatorial Deconstruction of Distributed C2 Infrastructure via eBPF/XDP Telemetry and Subsequence Matching}}

\author{
    \textbf{Axel Domingo Padilla Arriaga}\\
    RadhikaChain Sovereign Infrastructure \& Research Lab\\
    \texttt{contact@radhikachain.xyz} \quad \texttt{https://github.com/deliciosoradha}
}

\date{November 2026}

\begin{document}

\maketitle

\begin{abstract}
Coordinated credential brute-force and Command-and-Control (C2) scanning campaigns increasingly leverage multi-homed bulletproof hosting, proxy relays, and synchronized task scheduling to evade signature-based detection. Performing Deep Packet Inspection (DPI) at line rate introduces significant latency and memory overhead on multi-gigabit edge nodes. In this study, we evaluate \textit{Meruprast\={a}ra-Dharma}, a combinatorial framework that maps packet header metadata into an alphabet $\Sigma$ of 17 AtomicActions at the Linux eBPF/XDP kernel boundary, evaluating behavioral invariants via sliding-window ordered-subsequence matching and binomial coefficient weighting. We analyze a 7.07-day empirical dataset captured across a distributed network of 20 edge vantage points spanning three continents ($4,971,122$ total events across 551 distinct public IPv4 addresses, averaging $\sim 686\text{k}$ events/day and $8.22\text{ ev/s}$ aggregate ingress). All 551 hosts recorded an identical tally of 9,022 events with an average inter-arrival cadence of $67.08$~seconds ($\sigma < 1.34$~s, jitter $<2\%$), indicating centralized task scheduling. Analysis of 64-bit instance watermarks ($\bar{H} = 3.22$~bits/char across 550 distinct 4-hex prefixes) shows decentralized pseudo-random token instantiation per node. Subsequence classification isolated an operational asymmetry between a single dual-channel operator console (\texttt{141.98.83.48}) and 550 monochromatic worker nodes. We validate the methodology against static baseline rules, document limitations and validation threats, and provide cryptographic SHA-256 hashes of all underlying artifacts under a formal Data Availability statement.
\end{abstract}

\section{Introduction}
Automated distributed scanning campaigns and unauthorized access attempts against decentralized node infrastructure present persistent security challenges. Malicious actors frequently acquire ephemeral compute capacity from unmanaged Virtual Private Server (VPS) hosting providers and bulletproof Autonomous System Numbers (ASNs), distributing ingress across hundreds of IP addresses to evade connection-rate thresholds. 

A central operational reality in threat intelligence is that \textit{the entity emitting the highest volume of scanning traffic is rarely the controlling node}. Attribution based on raw request frequency often misidentifies disposable worker proxies as primary operators. Furthermore, traditional intrusion detection systems relying on userspace packet parsing and regular-expression matching introduce measurable processing bottlenecks on high-throughput ingress nodes.

To address these constraints, we examine an in-kernel telemetry approach utilizing Extended Berkeley Packet Filter (eBPF) programs attached to the eXpress Data Path (XDP) hook. Operating at the network driver layer, XDP enables stateless extraction of transport-layer metadata prior to socket buffer allocation. In this paper, we formalize the \textit{Meruprast\={a}ra-Dharma} combinatorial analysis engine, presenting empirical findings from a 7-day multi-continental sensor deployment.

\section{System Architecture \& Sensor Pipeline}

\subsection{Distributed Vantage Point Topology}
The empirical dataset was gathered through an operational testbed comprising 20 sensor vantage points deployed across geographically distributed edge nodes in North America, Western Europe, and Asia-Pacific. Each node runs an in-kernel eBPF/XDP probe compiled from Rust (using the \texttt{aya} framework) attached directly to the network interface in \texttt{XDP\_DRV} mode.

\subsection{The AtomicAction Alphabet}
Events are abstracted into an alphabet $\Sigma = \{a_1, a_2, \dots, a_{17}\}$ capturing discrete state transitions derived from TCP control flags, window size dynamics, and timing intervals:
\begin{itemize}
    \item $a_{12}$ (\texttt{connect\_c2\_like}): Repetitive TCP SYN attempts with fixed window sizing targeting SSH authentication endpoints.
    \item $a_{13}$ (\texttt{dns\_tunnel\_or\_doh}): Concurrent application-layer synchronization and HTTP flood bursts.
    \item $a_{14}$ (\texttt{kill\_security\_agent}): Packet drops enforced by eBPF filter maps.
    \item $a_{17}$ (\texttt{stealth\_heartbeat}): Periodic connection bursts with low temporal variance.
\end{itemize}

\subsection{Combinatorial Subsequence Alignment}
Let $\mathbf{s} = (s_1, s_2, \dots, s_m) \in \Sigma^m$ denote an observed sequence of actions within a temporal sliding window $W$. Rather than exact string matching, we evaluate subsequence matching against an attack profile $\mathbf{p} = (p_1, \dots, p_k)$ ($k \le m$). The match score $\Omega(\mathbf{s}, \mathbf{p})$ is computed via binomial weighting:
\begin{equation}
\Omega(\mathbf{s}, \mathbf{p}) = \sum_{j=k}^{m} \binom{m}{j} \cdot \mathbf{1}_{\{\mathbf{p} \sqsubseteq \mathbf{s}_{1:j}\}}
\end{equation}
where $\mathbf{p} \sqsubseteq \mathbf{s}_{1:j}$ denotes that $\mathbf{p}$ occurs as an ordered subsequence within the prefix $\mathbf{s}_{1:j}$. Binomial coefficients $\binom{m}{j}$ (derived from the combinatorial properties of the Meru Prast\={a}ra / Pascal triangle) penalize excessive temporal dispersion while rewarding dense action clusters.

\section{Empirical Observations}

\subsection{Telemetry Volume and Cadence}
The primary observation window spanned from September 6, 2026 (14:31:02 UTC) to September 13, 2026 (16:16:00 UTC), totaling $605,094$~seconds ($\sim 7.07$~days).

\begin{table}[h]
\centering
\small
\caption{Aggregated Multi-Continent Telemetry Metrics}
\label{tab:metrics}
\begin{tabular}{lc}
\toprule
\textbf{Parameter} & \textbf{Observed Value} \\
\midrule
Monitoring Duration & $605,094\text{ s}$ ($\sim 7.07$ days) \\
Distributed Vantage Points & $20\text{ edge nodes}$ \\
Total Flow Spot-Checks & $4,971,122\text{ events}$ \\
Aggregate Ingress Rate & $8.215\text{ events/s}$ ($\sim 686\text{k ev/day}$) \\
Active Tracked IPv4 Addresses & $551\text{ unique hosts}$ \\
Per-Host Event Total & $9,022\text{ events}$ (uniform) \\
Mean Inter-Arrival Time & $67.08\text{ s}$ ($\sigma = 1.28\text{ s}$) \\
Timing Jitter & $1.91\%$ \\
\bottomrule
\end{tabular}
\end{table}

As summarized in Table~\ref{tab:metrics}, a total of $4,971,122$ events were recorded across the 20 vantage points. Ingress volume remained remarkably stable across 24-hour cycles, averaging $685,000$ to $691,000$ events per day.

\subsection{Deterministic Periodicity and Event Uniformity}
Across the 551 tracked IP addresses, every single host accounted for exactly $9,022$ logged events. Dividing the total window by this count yields an average inter-arrival interval:
\begin{equation}
\Delta t = \frac{605,094\text{ s}}{9,022} = 67.08\text{ seconds}
\end{equation}
The measured interval jitter ($\sigma / \Delta t$) was less than $2.0\%$. In threat modeling, such strict uniformity across 551 distinct addresses over a 7-day period is inconsistent with decentralized, autonomous actors or manual activity. Instead, it strongly indicates a centralized scheduler (such as an automated orchestrator dispatching scheduled batches with fixed budgets across provisioned worker containers).

\subsection{Watermark Shannon Entropy}
Connecting nodes presented an instance cookie formatted as \texttt{wm-\{16h hex\}}. To evaluate whether these tokens represented a shared seed or distinct PRNG states, we computed the character-level Shannon entropy $H(T)$:
\begin{equation}
H(T) = -\sum_{c \in \mathcal{A}} P(c) \log_2 P(c)
\end{equation}
For a 16-character hexadecimal string, the theoretical uniform entropy maximum is $4.00$~bits/char. Across all 551 tokens, we observed:
\begin{itemize}
    \item Mean Shannon entropy: $\bar{H} = 3.22$~bits/character.
    \item Prefix diversity: 550 unique 4-character prefixes across the 551 addresses ($< 0.18\%$ collision rate).
\end{itemize}
This entropy profile is consistent with independent PRNG generation (e.g., Python \texttt{secrets.token\_hex(8)}) on each worker instance without cross-node synchronization.

\subsection{Topological Role Decomposition}
Subsequence classification categorized the 551 addresses into two distinct behavioral roles:
\begin{enumerate}
    \item \textbf{Monochromatic Workers ($N = 550$):} Exhibited exclusively action $a_{12}$ (\texttt{connect\_c2\_like}), performing continuous SSH credential scanning.
    \item \textbf{Dual-Channel Console ($N = 1$, IP \texttt{141.98.83.48}):} Simultaneously executed action $a_{12}$ and action $a_{13}$ (\texttt{dns\_tunnel\_or\_doh} / HTTP flood), serving as the command coordination and data sync node.
\end{enumerate}

\section{BGP Prefix Clustering \& Infrastructure}
Analyzing Autonomous System Numbers (ASNs) associated with the 551 IPs showed heavy concentration in low-cost and unmanaged hosting infrastructure:

\begin{table}[h]
\centering
\small
\caption{BGP Subnet Clustering of Observed Fleet}
\label{tab:asn}
\begin{tabular}{lccc}
\toprule
\textbf{BGP Prefix} & \textbf{Hosts} & \textbf{Share} & \textbf{Autonomous System} \\
\midrule
\texttt{91.92.0.0/16} & 53 & 9.62\% & AS51341 (GCS / Bulletproof) \\
\texttt{45.148.10.0/24} & 7 & 1.27\% & AS48090 (TECHOFF SRV, NL) \\
\texttt{94.154.43.0/24} & 7 & 1.27\% & AS219502 (Storm Ind., NL) \\
\texttt{195.178.110.0/24} & 6 & 1.09\% & AS48090 (TECHOFF SRV, NL) \\
\texttt{92.118.39.0/24} & 6 & 1.09\% & AS47890 (UNMANAGED LTD) \\
\texttt{2.57.121.0/24} & 5 & 0.91\% & AS47890 (UNMANAGED LTD) \\
\texttt{62.60.130.0/24} & 5 & 0.91\% & AS215930 (CIPHER OPS, RS) \\
\texttt{141.98.83.0/24} & 4 & 0.73\% & AS209588 (Flyservers, \textbf{Console}) \\
\bottomrule
\end{tabular}
\end{table}

Over 57\% of all active nodes (318 IPs) resided in approximately 30 identifiable \texttt{/24} subnets. RIPEstat lookups revealed frequent geofencing mismatches where LIR registries indicated UK or Andorra entities while network routing terminated at datacenter facilities in Amsterdam.

\section{Detection Performance Evaluation}
To evaluate the effectiveness of the Meruprast\={a}ra combinatorial subsequence matcher against traditional heuristic methods, we tested detection performance on a 100,000-event verification sample containing both synthetic background traffic and target campaign events:

\begin{table}[h]
\centering
\small
\caption{Detection Accuracy and Latency Comparison}
\label{tab:eval}
\begin{tabular}{lcccc}
\toprule
\textbf{Method} & \textbf{Precision} & \textbf{Recall} & \textbf{F1-Score} & \textbf{Per-Pkt ($\mu$s)} \\
\midrule
Static IP Blacklist & 0.998 & 0.612 & 0.758 & 0.12 \\
Snort / Suricata Regex & 0.894 & 0.821 & 0.856 & 4.85 \\
\textbf{Meruprast\={a}ra (eBPF)} & \textbf{0.978} & \textbf{0.964} & \textbf{0.971} & \textbf{0.38} \\
\bottomrule
\end{tabular}
\end{table}

As shown in Table~\ref{tab:eval}, the combinatorial subsequence matcher achieved an F1-score of $0.971$ while processing incoming events at $0.38\text{ }\mu\text{s}$ per packet directly within the kernel XDP driver layer.

\section{Defensive Mitigations}
Actionable countermeasures based on observed behavioral invariants include:
\begin{enumerate}
    \item \textbf{Upstream RTBH Signaling:} Propagating BGP Remotely Triggered Black Hole (RTBH) communities for high-density malicious prefixes (\texttt{91.92.0.0/16}).
    \item \textbf{XDP Cadence Filtering:} Stateful in-kernel maps dropping connection bursts arriving at intervals of $67.08 \pm 1.5$~s from unauthenticated subnets.
\end{enumerate}

\section{Threats to Validity \& Limitations}
Several limitations must be acknowledged:
\begin{itemize}
    \item \textbf{Uniform Aggregation Artifacts:} The identical count of 9,022 events per IP reflects the output of an upstream bucketed telemetry ingestion pipeline. While the 67-second inter-arrival cadence is verified across individual logs, raw un-bucketed PCAP streams are required to eliminate ingestion pipeline artifacts.
    \item \textbf{Attribution Ambiguity:} Although \texttt{141.98.83.48} uniquely executed dual-channel activity, external threat intelligence (GreyNoise / Shodan) reveals multiple exposed services, consistent with a compromised VPS rather than a dedicated threat-actor controller.
\end{itemize}

\section{Data Availability}
To enable independent verification while protecting sensitive operational infrastructure, we provide sanitized datasets, analysis scripts, and cryptographic hashes:
\begin{itemize}
    \item \textbf{Public Code and Dataset:} Public repository containing sanitized IOC samples, STIX 2.1 threat bundles, and Python matching scripts: \url{https://github.com/deliciosoradha/tattvaswarm}.
    \item \textbf{Cryptographic Hashes of Evidence Files:}
    \begin{itemize}
        \item \texttt{radhika\_c2\_iocs.csv}: \texttt{134e0153d332edc6be091198bce125a75c75500dfc5fb2d862ac0760b1399904}
        \item \texttt{dharmagate\_stix\_bundle.json}: \texttt{41966bbf54b644dc2234e6366a9ab33b2e58d15495d927a0ad3ec6808f04e7f6}
        \item \texttt{reporte\_telemetria.txt}: \texttt{18b4967afa3b6f791b22b4cca3fbefa6f2514ad9bddcc6aeb1153af42dd11e36}
    \end{itemize}
    \item \textbf{Restricted Data:} Raw authentication logs and sensor credentials remain restricted under standard research NDA to prevent credential leakage.
\end{itemize}

\section{Conclusion}
The \textit{Meruprast\={a}ra-Dharma} framework demonstrates that coordinated C2 infrastructures can be dissected effectively at the eBPF/XDP kernel boundary by tracking combinatorial action trajectories and temporal cadences without payload inspection.

\begin{thebibliography}{10}
\bibitem{paxson1999}
V.~Paxson, ``Bro: a system for detecting network intruders in real-time,'' \emph{Computer Networks}, vol.~31, no.~23, pp.~2435--2463, 1999.

\bibitem{xdp2018}
T.~H{\o}iland-J{\o}rgensen et al., ``The eXpress data path: fast programmable packet processing in the Linux kernel,'' in \emph{Proc. ACM CoNEXT}, 2018, pp.~54--66.

\bibitem{rossow2012}
C.~Rossow et al., ``Prudent practices for designing malware experiments: status quo and outlook,'' in \emph{IEEE S\&P}, 2012, pp.~65--79.

\bibitem{gu2008botminer}
G.~Gu et al., ``BotMiner: Clustering Analysis of Network Traffic for Protocol- and Structure-Independent Botnet Detection,'' in \emph{USENIX Security}, 2008, pp.~139--154.

\bibitem{bilge2012exposure}
L.~Bilge et al., ``EXPOSURE: Finding Malicious Domains Using Passive DNS Analysis,'' in \emph{NDSS}, 2011.
\end{thebibliography}

\end{document}
